% Generated by paper_mdpi/tools/make_tables.py. Do not edit by hand.
\begin{table}[H]
\caption{Speed and size of the operating points and of the transformer detectors\label{tab:operating}}
\begin{adjustwidth}{-\extralength}{0cm}
\small
\begin{tabularx}{\fulllength}{lcCCCCC}
\toprule
Model & Task & p50, ms & p95, ms & msg/s & MB & Routed, \% \\
\midrule
Lite (char SVM) & 3 classes & 1.47 & 1.87 & 660 & 5.4 & -- \\
Char SVM, aug. (stage 1 of Full) & 3 classes & 1.46 & 1.73 & 669 & 6.6 & -- \\
Full (cascade) & 3 classes & 1.60 & 2.47 & 297 & 272.9 & 2.1 \\
DistilBERT, aug. (stage 2 of Full) & 3 classes & 41.84 & 76.01 & 21 & 266.3 & -- \\
BERT & 3 classes & 84.67 & 159.07 & 10 & 438.7 & -- \\
\midrule
Shim et al., DistilBERT & Shim & 43.17 & 82.96 & 20 & 268.5 & -- \\
Shim et al., BERT & Shim & 87.20 & 161.21 & 10 & 438.7 & -- \\
Shim et al., RoBERTa & Shim & 89.53 & 170.75 & 10 & 502.2 & -- \\
Mahmud et al., CNN-BiGRU & Mahmud & 5.87 & 6.48 & 168 & 3.3 & -- \\
\bottomrule
\end{tabularx}
\noindent{\footnotesize{One CPU thread, batch size 1, seed-1 weights; task: the test messages used for timing. Full is measured end to end on the whole three-class test set, including the calls of its second stage; Routed: share of the clean test messages sent to the second stage, mean over 5 seeds. Size: saved model files. All models of the three-class task: Table~\ref{tab:efficiency}.}}
\end{adjustwidth}
\end{table}
